% Barebone template for use with Zettlr exporting engine
% We have to use KOMA classes (scrartcl instead of article and scrreprt instead of report as well as scrbook instead of book)
% More info: https://ctan.org/pkg/koma-script?lang=de
\documentclass[
$FONT_SIZE$,
$if(lang)$
  $babel-lang$,
$endif$
]{scrartcl}

\newcommand{\documentAuthor}{$PDF_AUTHOR$}
\newcommand{\documentTitle}{$PDF_TITLE$}

%%%%SETTINGS%%%%
%% Document language (en, de)
\newcommand{\documentLanguage}{de}
\input{ads/lang/\documentLanguage.tex}

%% Document type
% T2\_1000 Project Thesis (Semester 1 & 2)
% T2\_2000 Project Thesis (Semester 3 & 4)
% T2\_3100 Seminar Paper (Semester 5 & 6)
% T2\_3300 Bachelor Thesis
\newcommand{\documentType}{T2\_3000}
%\newcommand{\documentTypePhrase}{Projektarbeit}

\newcommand{\multipleAuthors}{false}

\newcommand{\documentPeriod}{12 Wochen}

\newcommand{\matriculationNumber}{4303139}

\newcommand{\locationUniversity}{Friedrichshafen}
\newcommand{\department}{Wirtschatsingenieurwesen - Maschinenbau}
\newcommand{\course}{TWI18-1}

\newcommand{\degree}{Bachelor of Engineering}
% INF2014 - INF2016 (MI):       Bachelor of Science
% INF2014 - INF2016 (IA/IM) :   Bachelor of Engineering
% INF2017 (all):                Bachelor of Science

\newcommand{\releaseDate}{März 2021}
\newcommand{\releaseLocation}{Manching}

\newcommand{\companyName}{Airbus Defence and Space GmbH}
\newcommand{\companyLocation}{Manching}

\newcommand{\tutor}{B.Sc. Andreas Keis}
\newcommand{\evaluator}{}

\newcommand{\linkColor}{505050}
%%%%SETTINGS%%%%%
%%%%HEADER%%%%
\usepackage{xstring}
\usepackage[utf8]{inputenc}
\usepackage[T1]{fontenc}


% iflang command definition
\newcommand{\iflang}[2]{%
    \IfStrEq{\documentLanguage}{#1}{#2}{}%
}

% ifDocType comand definition
\newcommand{\ifDocType}[3]{%
  \IfStrEq{\documentType}{#1}{#2}{#3}%
}

% ifMultipleAuthors definition
\newcommand{\ifMultipleAuthors}[2]{%
  \IfStrEq{\multipleAuthors}{true}{#1}{#2}%
}

% Include main settings
%% Barebone template for use with Zettlr exporting engine
% We have to use KOMA classes (scrartcl instead of article and scrreprt instead of report as well as scrbook instead of book)
% More info: https://ctan.org/pkg/koma-script?lang=de
\documentclass[
$FONT_SIZE$,
$if(lang)$
  $babel-lang$,
$endif$
]{scrartcl}

\newcommand{\documentAuthor}{$PDF_AUTHOR$}
\newcommand{\documentTitle}{$PDF_TITLE$}

%%%%SETTINGS%%%%
%% Document language (en, de)
\newcommand{\documentLanguage}{de}
\input{ads/lang/\documentLanguage.tex}

%% Document type
% T2\_1000 Project Thesis (Semester 1 & 2)
% T2\_2000 Project Thesis (Semester 3 & 4)
% T2\_3100 Seminar Paper (Semester 5 & 6)
% T2\_3300 Bachelor Thesis
\newcommand{\documentType}{T2\_3000}
%\newcommand{\documentTypePhrase}{Projektarbeit}

\newcommand{\multipleAuthors}{false}

\newcommand{\documentPeriod}{12 Wochen}

\newcommand{\matriculationNumber}{4303139}

\newcommand{\locationUniversity}{Friedrichshafen}
\newcommand{\department}{Wirtschatsingenieurwesen - Maschinenbau}
\newcommand{\course}{TWI18-1}

\newcommand{\degree}{Bachelor of Engineering}
% INF2014 - INF2016 (MI):       Bachelor of Science
% INF2014 - INF2016 (IA/IM) :   Bachelor of Engineering
% INF2017 (all):                Bachelor of Science

\newcommand{\releaseDate}{März 2021}
\newcommand{\releaseLocation}{Manching}

\newcommand{\companyName}{Airbus Defence and Space GmbH}
\newcommand{\companyLocation}{Manching}

\newcommand{\tutor}{B.Sc. Andreas Keis}
\newcommand{\evaluator}{}

\newcommand{\linkColor}{505050}
%%%%SETTINGS%%%%%
%%%%HEADER%%%%
\usepackage{xstring}
\usepackage[utf8]{inputenc}
\usepackage[T1]{fontenc}


% iflang command definition
\newcommand{\iflang}[2]{%
    \IfStrEq{\documentLanguage}{#1}{#2}{}%
}

% ifDocType comand definition
\newcommand{\ifDocType}[3]{%
  \IfStrEq{\documentType}{#1}{#2}{#3}%
}

% ifMultipleAuthors definition
\newcommand{\ifMultipleAuthors}[2]{%
  \IfStrEq{\multipleAuthors}{true}{#1}{#2}%
}

% Include main settings
%% Barebone template for use with Zettlr exporting engine
% We have to use KOMA classes (scrartcl instead of article and scrreprt instead of report as well as scrbook instead of book)
% More info: https://ctan.org/pkg/koma-script?lang=de
\documentclass[
$FONT_SIZE$,
$if(lang)$
  $babel-lang$,
$endif$
]{scrartcl}

\newcommand{\documentAuthor}{$PDF_AUTHOR$}
\newcommand{\documentTitle}{$PDF_TITLE$}

%%%%SETTINGS%%%%
%% Document language (en, de)
\newcommand{\documentLanguage}{de}
\input{ads/lang/\documentLanguage.tex}

%% Document type
% T2\_1000 Project Thesis (Semester 1 & 2)
% T2\_2000 Project Thesis (Semester 3 & 4)
% T2\_3100 Seminar Paper (Semester 5 & 6)
% T2\_3300 Bachelor Thesis
\newcommand{\documentType}{T2\_3000}
%\newcommand{\documentTypePhrase}{Projektarbeit}

\newcommand{\multipleAuthors}{false}

\newcommand{\documentPeriod}{12 Wochen}

\newcommand{\matriculationNumber}{4303139}

\newcommand{\locationUniversity}{Friedrichshafen}
\newcommand{\department}{Wirtschatsingenieurwesen - Maschinenbau}
\newcommand{\course}{TWI18-1}

\newcommand{\degree}{Bachelor of Engineering}
% INF2014 - INF2016 (MI):       Bachelor of Science
% INF2014 - INF2016 (IA/IM) :   Bachelor of Engineering
% INF2017 (all):                Bachelor of Science

\newcommand{\releaseDate}{März 2021}
\newcommand{\releaseLocation}{Manching}

\newcommand{\companyName}{Airbus Defence and Space GmbH}
\newcommand{\companyLocation}{Manching}

\newcommand{\tutor}{B.Sc. Andreas Keis}
\newcommand{\evaluator}{}

\newcommand{\linkColor}{505050}
%%%%SETTINGS%%%%%
%%%%HEADER%%%%
\usepackage{xstring}
\usepackage[utf8]{inputenc}
\usepackage[T1]{fontenc}


% iflang command definition
\newcommand{\iflang}[2]{%
    \IfStrEq{\documentLanguage}{#1}{#2}{}%
}

% ifDocType comand definition
\newcommand{\ifDocType}[3]{%
  \IfStrEq{\documentType}{#1}{#2}{#3}%
}

% ifMultipleAuthors definition
\newcommand{\ifMultipleAuthors}[2]{%
  \IfStrEq{\multipleAuthors}{true}{#1}{#2}%
}

% Include main settings
%% Barebone template for use with Zettlr exporting engine
% We have to use KOMA classes (scrartcl instead of article and scrreprt instead of report as well as scrbook instead of book)
% More info: https://ctan.org/pkg/koma-script?lang=de
\documentclass[
$FONT_SIZE$,
$if(lang)$
  $babel-lang$,
$endif$
]{scrartcl}

\newcommand{\documentAuthor}{$PDF_AUTHOR$}
\newcommand{\documentTitle}{$PDF_TITLE$}

\input{ads/settings.tex}

% iflang command definition
\newcommand{\iflang}[2]{%
    \IfStrEq{\documentLanguage}{#1}{#2}{}%
}

% ifDocType comand definition
\newcommand{\ifDocType}[3]{%
  \IfStrEq{\documentType}{#1}{#2}{#3}%
}

% ifMultipleAuthors definition
\newcommand{\ifMultipleAuthors}[2]{%
  \IfStrEq{\multipleAuthors}{true}{#1}{#2}%
}

% Include main settings
%\input{settings/main} 

% Include document settings
%\input{settings/document}

% Load language specific Strings
%\input{lang/en.tex}

% Add comment feature
\newcommand{\comment}[1]{\par {\bfseries \color{blue} #1 \par}}


%%%%%%% Package Includes %%%%%%%
\usepackage{chngcntr}
\usepackage{float}
\usepackage{subfig}
\usepackage{tabularx}
%\usepackage[margin=\margin,foot=1cm]{geometry}
%\usepackage[activate]{microtype}                                         
%\usepackage[onehalfspacing]{setspace}
\usepackage{makeidx}
\usepackage[autostyle=true,german=quotes]{csquotes}
\usepackage{longtable}
\usepackage{enumitem}	                                                 
\usepackage{graphicx}
\usepackage{pdfpages}                                                         
\usepackage{xcolor} 	                                                    
\usepackage{float}
\usepackage{array}
\usepackage{calc}		                         
\usepackage[right]{eurosym}
\usepackage{wrapfig}
\usepackage{pgffor}                                                             
\usepackage[perpage, hang, multiple, stable]{footmisc}  
\usepackage[printonlyused]{acronym}                                 
\usepackage{listings}
\usepackage[obeyFinal,backgroundcolor=yellow,linecolor=black]{todonotes}
\usepackage{rotating}
\usepackage{lscape}
\usepackage{amsmath}
\usepackage{amssymb}
\usepackage{eurosym}
\usepackage{tocbibind} %makes list of figures, list of tables etc. show up in toc
%\usepackage{\documentFont}
%\usepackage[
%	pdftitle={\documentTitle},
%	pdfauthor={\documentAuthor},
%	pdfsubject={\documentType},
%	pdfcreator={pdflatex, LaTeX with KOMA-Script},
%	pdfpagemode=UseOutlines,       % Show Contents while opening
	%pdfdisplaydoctitle=true, 		% Show document title instead of file name
	%pdflang={\documentLanguage}, % Document language
%]{hyperref}
\usepackage{bookmark}
\usepackage[nonumberlist,toc,xindy]{glossaries}
% add first occurence in footnote
\defglsdisplayfirst[main]{#1#4\protect\footnote{#2}}

% Commands for check and x marks
\usepackage{amssymb}% http://ctan.org/pkg/amssymb
\usepackage{pifont}% http://ctan.org/pkg/pifont
\newcommand{\cmark}{\ding{51}}%
\newcommand{\xmark}{\ding{55}}%
    
% Generate glossary
\makeglossaries{}

% Load colors
%\defineColors{}

% Set Titel, Autor and Date
%\title{\documentTitle}
%\author{\documentAuthor}
%\date{\datum}


% PDF link settings
\hypersetup{%
	colorlinks=true, 		
	linkcolor=LinkColor, 	
	citecolor=LinkColor,
	filecolor=LinkColor,
	menucolor=LinkColor,
	urlcolor=LinkColor,
	linktocpage=true, 
	bookmarksnumbered=true 
}

% Captions fontsize
\addtokomafont{caption}{\small}

% Hurenkinder und Schusterjungen verhindern
% http://projekte.dante.de/DanteFAQ/Silbentrennung
\clubpenalty = 10000 % schließt Schusterjungen aus (Seitenumbruch nach der ersten Zeile eines neuen Absatzes)
\widowpenalty = 10000 % schließt Hurenkinder aus (die letzte Zeile eines Absatzes steht auf einer neuen Seite)
\displaywidowpenalty=10000

% Graphicspath
\graphicspath{{images/}}

% frequently used programing languages
\lstloadlanguages{PHP,Python,Java,C,C++,bash}

%\listingsettings{}
% Rename Listings
\renewcommand\lstlistingname{\listingPhrase}
\renewcommand\lstlistlistingname{\listListingPhrase}
\def\lstlistingautorefname{\authorListingPhrase}

% Spaces in tables
%\setlength{\tabcolsep}{\tableColumnMargin}
%\renewcommand{\arraystretch}{\tableRowMargin}

% Manage ToDo's
%\usepackage{todonotes}
%\usepackage{ulem}
\iflang{de}{\usepackage[english, ngerman]{babel}}
\iflang{en}{\usepackage[ngerman, english]{babel}}
%%%%HEADER%%%%

% The following packages are needed for the combined use of Pandoc and XeTeX engine
\usepackage{fontspec,xltxtra,xunicode}

% These directives are necessary for Pandoc to correctly set the language of the
% document. If the user uses XeLaTeX (the default for Zettlr), it will use
% polyglossia, else it will use babel. We've simply copied that from the default
% template -- thanks to the Pandoc community at this point!
$if(lang)$
\ifxetex
  \usepackage{polyglossia}
  \setmainlanguage[$polyglossia-lang.options$]{$polyglossia-lang.name$}
$for(polyglossia-otherlangs)$
  \setotherlanguage[$polyglossia-otherlangs.options$]{$polyglossia-otherlangs.name$}
$endfor$
\else
  \usepackage[shorthands=off,$for(babel-otherlangs)$$babel-otherlangs$,$endfor$main=$babel-lang$]{babel}
\fi
$endif$

% Captions for tables and images
\usepackage{caption}
\captionsetup{format=plain, font=small, labelfont=bf}

% This package is needed for better rendering of images (including captions)
\usepackage{graphicx}

% These packages are needed by pandoc for certain special glyphs, such as the
% \square symbol for checkboxes.
\usepackage{amssymb,amsmath}

% Calc is needed to compute the image width (to prevent up-scaling of small images)
\usepackage{calc}

% ulem is needed for underlining and strikethrough text commands.
% the option "normalem" is needed to preserve _italics_. Without the option,
% _italics_ will become underlines, which is not desirable.
\usepackage[normalem]{ulem}

% Image scaling code from pandoc default.latex template
\makeatletter
\def\maxwidth{\ifdim\Gin@nat@width>\linewidth\linewidth\else\Gin@nat@width\fi}
\def\maxheight{\ifdim\Gin@nat@height>\textheight\textheight\else\Gin@nat@height\fi}
\makeatother
% Scale images if necessary, so that they will not overflow the page
% margins by default, and it is still possible to overwrite the defaults
% using explicit options in \includegraphics[width, height, ...]{}
\setkeys{Gin}{width=\maxwidth,height=\maxheight,keepaspectratio}

% Now we also have to make sure that the figures Pandoc inserts are positioned
% approximately where they are at. Therefore we have to overwrite the
% \begin{figure} command a little bit by using the float package:
\usepackage{float}
\floatplacement{figure}{H} % ensure default position as H (exactly where the figure is in the text)

% Needed for pandoc's table generation
\usepackage{longtable, booktabs, array}

% Optional packages that Zettlr makes use of
% fontspec is used to set the different fonts to the user's whishes
\setmainfont{$MAIN_FONT$}
\setsansfont{$SANS_FONT$}

% With geometry we can let the user decide on the margin
\usepackage{geometry}
% Additional options for geometry can be found here: http://texdoc.net/texmf-dist/doc/latex/geometry/geometry.pdf
\geometry{verbose,$PAPER_TYPE$,tmargin=$TOP_MARGIN$,bmargin=$BOTTOM_MARGIN$,lmargin=$LEFT_MARGIN$,rmargin=$RIGHT_MARGIN$}

% PAGE NUMBERING OPTIONS: https://de.sharelatex.com/learn/Page_numbering
\pagenumbering{$PAGE_NUMBERING$}

% Line spacing
\usepackage{setspace}

% Can be a multiple of 1 (i.e. 1.5 for 150 percent)
\setstretch{$LINE_SPACING$}

% PDF metadata generated by Zettlr or given by the user, also some other URL
% options (breaklinks breaks URLs so they aren't ugly)
%\usepackage[hyphens]{url} % This package breaks links even better
\usepackage{hyperref}
\hypersetup{
unicode=true,
breaklinks=true,
colorlinks=true,
linkcolor=[rgb]{0.00,0.00,0.00},
urlcolor=[rgb]{0.00,0.00,0.00},
pdftitle={$PDF_TITLE$},
pdfsubject={$PDF_SUBJECT$},
pdfauthor={$PDF_AUTHOR$},
pdfkeywords={$PDF_KEYWORDS$},
pdfproducer={Zettlr with Pandoc and XeLaTeX},
pdfborder={0 0 0}
}

% Setup pandoc code highlighting
% pandoc uses an if statement, but until we have code
% to enable or disable highlighting we'll ignore that.

$highlighting-macros$


% Set the date (either \today or a date from the frontmatter)
\date{$PDF_DATE$}

% The following commands must be provided for Pandoc to work correctly.
\providecommand{\tightlist}{%
  \setlength{\itemsep}{0pt}\setlength{\parskip}{0pt}}

% The following commands are used by Pandoc if a CSL bibliography is present
% to create proper indentation, especially for styles with hanging indentation.
$if(csl-refs)$
\newlength{\cslhangindent}
\setlength{\cslhangindent}{1.5em}
\newlength{\csllabelwidth}
\setlength{\csllabelwidth}{3em}
\newenvironment{CSLReferences}[2] % #1 hanging-ident, #2 entry spacing
 {% don't indent paragraphs
  \setlength{\parindent}{0pt}
  % turn on hanging indent if param 1 is 1
  \ifodd #1 \everypar{\setlength{\hangindent}{\cslhangindent}}\ignorespaces\fi
  % set entry spacing
  \ifnum #2 > 0
  \setlength{\parskip}{#2\baselineskip}
  \fi
 }%
 {}
\usepackage{calc}
\newcommand{\CSLBlock}[1]{#1\hfill\break}
\newcommand{\CSLLeftMargin}[1]{\parbox[t]{\csllabelwidth}{#1}}
\newcommand{\CSLRightInline}[1]{\parbox[t]{\linewidth - \csllabelwidth}{#1}\break}
\newcommand{\CSLIndent}[1]{\hspace{\cslhangindent}#1}
$endif$

\title{$PDF_TITLE$}
\author{$PDF_AUTHOR$}

%\input{/home/damion/Desktop/Studium/Studienarbeit/Export/settings/document.tex}
%\input{/home/damion/Desktop/Studium/Studienarbeit/Export/ads/header.tex}

\begin{document}


% Title page?
%$TITLEPAGE$
\pagenumbering{Roman}
\input{ads/cover.tex}
\input{ads/declaration.tex}
\input{ads/restrictionNotices.tex}
\input{ads/acronyms.tex}

% Generate a table of contents?
$GENERATE_TOC$

% BEGIN BODY
\pagenumbering{$PAGE_NUMBERING$}
\setcounter{page}{1}
\setcounter{table}{0}
$body$

\pagenumbering{Roman}
\setcounter{page}{6}
\pagebreak
\listoftables
\pagebreak
\listoffigures
\end{document} 
 

% Include document settings
%\input{settings/document}

% Load language specific Strings
%%!TEX root = ../main.tex

\newcommand{\degreePhrase}{for the}
\newcommand{\departmentPhrase}{from the Course of Studies}
\newcommand{\locationUniversityPhrase}{at the Cooperative State University Baden-Württemberg}
\newcommand{\documentAuthorPhrase}{by}
\newcommand{\documentPeriodPhrase}{Time of Project}
\newcommand{\matriculationNumberPhrase}{Student ID}
\newcommand{\coursePhrase}{Course}
\newcommand{\companyPhrase}{Company}

\newcommand{\tutorPhrase}{%
	\IfStrEqCase{\documentType}{%
		    {T2\_1000}{Reviewer}%
      	{T2\_2000}{Reviewer}%
        {T2\_3000}{Reviewer}%
      	{T2\_3100}{Reviewer}%
      	{T2\_3300}{First Reviewer}%
	}
}

\newcommand{\evaluatorPhrase}{Second Reviewer}
\newcommand{\restrictionNoticesPhrase}{Confidentiality Statement} %Restriction Notice}
\newcommand{\declarationPhrase}{Author's declaration}
\newcommand{\glossaryPhrase}{Glossary}
\newcommand{\langabkverz}{Acronyms}
\newcommand{\langabstract}{Abstract}
\newcommand{\appendixPhrase}{Appendix}
\newcommand{\listingPhrase}{Listing}
\newcommand{\listListingPhrase}{Listings}
\newcommand{\authorListingPhrase}{Listing}
\newcommand{\acronymsPhrase}{Acronyms}

\newcommand{\documentTypePhrase}{%
  \IfStrEqCase{\documentType}{%
      {T2\_1000}{Project Thesis}%
      {T2\_2000}{Project Thesis}%
      {T2\_3000}{Project Thesis}%
      {T2\_3100}{Semester Paper}%
      {T2\_3300}{Bachelor Thesis}%
  }
}

% Add comment feature
\newcommand{\comment}[1]{\par {\bfseries \color{blue} #1 \par}}


%%%%%%% Package Includes %%%%%%%
\usepackage{chngcntr}
\usepackage{float}
\usepackage{subfig}
\usepackage{tabularx}
%\usepackage[margin=\margin,foot=1cm]{geometry}
%\usepackage[activate]{microtype}                                         
%\usepackage[onehalfspacing]{setspace}
\usepackage{makeidx}
\usepackage[autostyle=true,german=quotes]{csquotes}
\usepackage{longtable}
\usepackage{enumitem}	                                                 
\usepackage{graphicx}
\usepackage{pdfpages}                                                         
\usepackage{xcolor} 	                                                    
\usepackage{float}
\usepackage{array}
\usepackage{calc}		                         
\usepackage[right]{eurosym}
\usepackage{wrapfig}
\usepackage{pgffor}                                                             
\usepackage[perpage, hang, multiple, stable]{footmisc}  
\usepackage[printonlyused]{acronym}                                 
\usepackage{listings}
\usepackage[obeyFinal,backgroundcolor=yellow,linecolor=black]{todonotes}
\usepackage{rotating}
\usepackage{lscape}
\usepackage{amsmath}
\usepackage{amssymb}
\usepackage{eurosym}
\usepackage{tocbibind} %makes list of figures, list of tables etc. show up in toc
%\usepackage{\documentFont}
%\usepackage[
%	pdftitle={\documentTitle},
%	pdfauthor={\documentAuthor},
%	pdfsubject={\documentType},
%	pdfcreator={pdflatex, LaTeX with KOMA-Script},
%	pdfpagemode=UseOutlines,       % Show Contents while opening
	%pdfdisplaydoctitle=true, 		% Show document title instead of file name
	%pdflang={\documentLanguage}, % Document language
%]{hyperref}
\usepackage{bookmark}
\usepackage[nonumberlist,toc,xindy]{glossaries}
% add first occurence in footnote
\defglsdisplayfirst[main]{#1#4\protect\footnote{#2}}

% Commands for check and x marks
\usepackage{amssymb}% http://ctan.org/pkg/amssymb
\usepackage{pifont}% http://ctan.org/pkg/pifont
\newcommand{\cmark}{\ding{51}}%
\newcommand{\xmark}{\ding{55}}%
    
% Generate glossary
\makeglossaries{}

% Load colors
%\defineColors{}

% Set Titel, Autor and Date
%\title{\documentTitle}
%\author{\documentAuthor}
%\date{\datum}


% PDF link settings
\hypersetup{%
	colorlinks=true, 		
	linkcolor=LinkColor, 	
	citecolor=LinkColor,
	filecolor=LinkColor,
	menucolor=LinkColor,
	urlcolor=LinkColor,
	linktocpage=true, 
	bookmarksnumbered=true 
}

% Captions fontsize
\addtokomafont{caption}{\small}

% Hurenkinder und Schusterjungen verhindern
% http://projekte.dante.de/DanteFAQ/Silbentrennung
\clubpenalty = 10000 % schließt Schusterjungen aus (Seitenumbruch nach der ersten Zeile eines neuen Absatzes)
\widowpenalty = 10000 % schließt Hurenkinder aus (die letzte Zeile eines Absatzes steht auf einer neuen Seite)
\displaywidowpenalty=10000

% Graphicspath
\graphicspath{{images/}}

% frequently used programing languages
\lstloadlanguages{PHP,Python,Java,C,C++,bash}

%\listingsettings{}
% Rename Listings
\renewcommand\lstlistingname{\listingPhrase}
\renewcommand\lstlistlistingname{\listListingPhrase}
\def\lstlistingautorefname{\authorListingPhrase}

% Spaces in tables
%\setlength{\tabcolsep}{\tableColumnMargin}
%\renewcommand{\arraystretch}{\tableRowMargin}

% Manage ToDo's
%\usepackage{todonotes}
%\usepackage{ulem}
\iflang{de}{\usepackage[english, ngerman]{babel}}
\iflang{en}{\usepackage[ngerman, english]{babel}}
%%%%HEADER%%%%

% The following packages are needed for the combined use of Pandoc and XeTeX engine
\usepackage{fontspec,xltxtra,xunicode}

% These directives are necessary for Pandoc to correctly set the language of the
% document. If the user uses XeLaTeX (the default for Zettlr), it will use
% polyglossia, else it will use babel. We've simply copied that from the default
% template -- thanks to the Pandoc community at this point!
$if(lang)$
\ifxetex
  \usepackage{polyglossia}
  \setmainlanguage[$polyglossia-lang.options$]{$polyglossia-lang.name$}
$for(polyglossia-otherlangs)$
  \setotherlanguage[$polyglossia-otherlangs.options$]{$polyglossia-otherlangs.name$}
$endfor$
\else
  \usepackage[shorthands=off,$for(babel-otherlangs)$$babel-otherlangs$,$endfor$main=$babel-lang$]{babel}
\fi
$endif$

% Captions for tables and images
\usepackage{caption}
\captionsetup{format=plain, font=small, labelfont=bf}

% This package is needed for better rendering of images (including captions)
\usepackage{graphicx}

% These packages are needed by pandoc for certain special glyphs, such as the
% \square symbol for checkboxes.
\usepackage{amssymb,amsmath}

% Calc is needed to compute the image width (to prevent up-scaling of small images)
\usepackage{calc}

% ulem is needed for underlining and strikethrough text commands.
% the option "normalem" is needed to preserve _italics_. Without the option,
% _italics_ will become underlines, which is not desirable.
\usepackage[normalem]{ulem}

% Image scaling code from pandoc default.latex template
\makeatletter
\def\maxwidth{\ifdim\Gin@nat@width>\linewidth\linewidth\else\Gin@nat@width\fi}
\def\maxheight{\ifdim\Gin@nat@height>\textheight\textheight\else\Gin@nat@height\fi}
\makeatother
% Scale images if necessary, so that they will not overflow the page
% margins by default, and it is still possible to overwrite the defaults
% using explicit options in \includegraphics[width, height, ...]{}
\setkeys{Gin}{width=\maxwidth,height=\maxheight,keepaspectratio}

% Now we also have to make sure that the figures Pandoc inserts are positioned
% approximately where they are at. Therefore we have to overwrite the
% \begin{figure} command a little bit by using the float package:
\usepackage{float}
\floatplacement{figure}{H} % ensure default position as H (exactly where the figure is in the text)

% Needed for pandoc's table generation
\usepackage{longtable, booktabs, array}

% Optional packages that Zettlr makes use of
% fontspec is used to set the different fonts to the user's whishes
\setmainfont{$MAIN_FONT$}
\setsansfont{$SANS_FONT$}

% With geometry we can let the user decide on the margin
\usepackage{geometry}
% Additional options for geometry can be found here: http://texdoc.net/texmf-dist/doc/latex/geometry/geometry.pdf
\geometry{verbose,$PAPER_TYPE$,tmargin=$TOP_MARGIN$,bmargin=$BOTTOM_MARGIN$,lmargin=$LEFT_MARGIN$,rmargin=$RIGHT_MARGIN$}

% PAGE NUMBERING OPTIONS: https://de.sharelatex.com/learn/Page_numbering
\pagenumbering{$PAGE_NUMBERING$}

% Line spacing
\usepackage{setspace}

% Can be a multiple of 1 (i.e. 1.5 for 150 percent)
\setstretch{$LINE_SPACING$}

% PDF metadata generated by Zettlr or given by the user, also some other URL
% options (breaklinks breaks URLs so they aren't ugly)
%\usepackage[hyphens]{url} % This package breaks links even better
\usepackage{hyperref}
\hypersetup{
unicode=true,
breaklinks=true,
colorlinks=true,
linkcolor=[rgb]{0.00,0.00,0.00},
urlcolor=[rgb]{0.00,0.00,0.00},
pdftitle={$PDF_TITLE$},
pdfsubject={$PDF_SUBJECT$},
pdfauthor={$PDF_AUTHOR$},
pdfkeywords={$PDF_KEYWORDS$},
pdfproducer={Zettlr with Pandoc and XeLaTeX},
pdfborder={0 0 0}
}

% Setup pandoc code highlighting
% pandoc uses an if statement, but until we have code
% to enable or disable highlighting we'll ignore that.

$highlighting-macros$


% Set the date (either \today or a date from the frontmatter)
\date{$PDF_DATE$}

% The following commands must be provided for Pandoc to work correctly.
\providecommand{\tightlist}{%
  \setlength{\itemsep}{0pt}\setlength{\parskip}{0pt}}

% The following commands are used by Pandoc if a CSL bibliography is present
% to create proper indentation, especially for styles with hanging indentation.
$if(csl-refs)$
\newlength{\cslhangindent}
\setlength{\cslhangindent}{1.5em}
\newlength{\csllabelwidth}
\setlength{\csllabelwidth}{3em}
\newenvironment{CSLReferences}[2] % #1 hanging-ident, #2 entry spacing
 {% don't indent paragraphs
  \setlength{\parindent}{0pt}
  % turn on hanging indent if param 1 is 1
  \ifodd #1 \everypar{\setlength{\hangindent}{\cslhangindent}}\ignorespaces\fi
  % set entry spacing
  \ifnum #2 > 0
  \setlength{\parskip}{#2\baselineskip}
  \fi
 }%
 {}
\usepackage{calc}
\newcommand{\CSLBlock}[1]{#1\hfill\break}
\newcommand{\CSLLeftMargin}[1]{\parbox[t]{\csllabelwidth}{#1}}
\newcommand{\CSLRightInline}[1]{\parbox[t]{\linewidth - \csllabelwidth}{#1}\break}
\newcommand{\CSLIndent}[1]{\hspace{\cslhangindent}#1}
$endif$

\title{$PDF_TITLE$}
\author{$PDF_AUTHOR$}

%\input{/home/damion/Desktop/Studium/Studienarbeit/Export/settings/document.tex}
%\input{/home/damion/Desktop/Studium/Studienarbeit/Export/ads/header.tex}

\begin{document}


% Title page?
%$TITLEPAGE$
\pagenumbering{Roman}
%!TEX root = ../main.tex

\begin{titlepage}
	\begin{longtable}{p{8cm} p{8cm}}
		\raggedright {\raisebox{\ht\strutbox-\totalheight}{\includegraphics[height=2.5cm]{images/cover/logo-dhbw.png}}} &
		\ifDocType{T2\_3100}{%
			% no company logo for semester paper
		}{%
			\raggedleft {\raisebox{\ht\strutbox-\totalheight}{\includegraphics[height=2.5cm]{images/cover/logo-company.png}}}
		}%
	\end{longtable}
	\enlargethispage{20mm}
	\begin{center}
	\doublespacing{
		\vspace*{12mm}	{\LARGE\textbf \documentTitle }}\\
		\vspace*{12mm}	{\large\textbf {\documentTypePhrase}}\\
		
		% degre only for bachelor thesis
		\ifDocType{T2\_3300}{
			\vspace*{12mm}	\degreePhrase\\
			\vspace*{3mm}		{\textbf \degree}\\
		}

		\vspace*{12mm}	\departmentPhrase{} \department\\
	\vspace*{0mm}		\locationUniversityPhrase{} \locationUniversity\\
		\vspace*{12mm}	\documentAuthorPhrase\\
		\vspace*{3mm}		{\large\textbf \documentAuthor}\\
		\vspace*{12mm}	\releaseDate\\
	\end{center}
	\vfill
	\begin{spacing}{1.2}
	\begin{tabbing}
		mmmmmmmmmmmmmmmmmmmmmmmmmm             \= \kill
		\textbf{\documentPeriodPhrase}       \>  \documentPeriod\\
		\textbf{\matriculationNumberPhrase, \coursePhrase}  \>  \matriculationNumber, \course\\

		\ifDocType{T2\_3100}{%
			% no company for the semester paper
		}{%
			\textbf{\companyPhrase}                  \>  \companyName, \companyLocation\\
		}%
		\textbf{\tutorPhrase}               \>  \tutor\\

		% evaluator only for bachelor thesis
		\ifDocType{T2\_3300}{%
			\textbf{\evaluatorPhrase}              \>  \evaluator
		}{}
	\end{tabbing}
	\end{spacing}
\end{titlepage}

%\thispagestyle{empty}

\section*{Erklärung}

\vspace*{2em}

gemäß Ziffer 1.1.13 der Anlage 1 zu §§ 3, 4 und 5  der Studien- und Prüfungsordnung für die Bachelorstudiengänge im Studienbereich Technik der Dualen Hochschule Baden-Württemberg vom 29.09.2017.
  
\ifMultipleAuthors{Wir versichern}{Ich versichere} hiermit, dass \ifMultipleAuthors{wir unsere}{ich meine} 
{\documentTypePhrase} mit dem Thema: 
\vspace{3em}
\begin{addmargin}[25pt]{0pt}
\Large {\itshape \documentTitle }
\end{addmargin}
\vspace{3em}
selbstständig verfasst und  keine anderen als die angegebenen Quellen und Hilfsmittel benutzt \ifMultipleAuthors{haben}{habe}. 
\ifMultipleAuthors{Wir versichern}{Ich versichere} zudem, dass die eingereichte elektronische Fassung mit der gedruckten Fassung 
übereinstimmt. 

\vspace{3em}

\releaseLocation, \releaseDate
\vspace{4em}

\rule{6cm}{0.4pt}

\documentAuthor
\pagebreak

%!TEX root = ../main.tex

%\thispagestyle{empty}

\section*{Sperrvermerk}

\vspace*{2em}

Der Inhalt dieser Arbeit darf weder als Ganzes noch in Auszügen Personen außerhalb des Prüfungsprozesses und des Evaluationsverfahrens zugänglich gemacht werden, sofern keine anderslautende Genehmigung vom Dualen Partner vorliegt.

\vspace{3em}

\releaseLocation, \releaseDate
\vspace{4em}

\rule{6cm}{0.4pt}

\documentAuthor
\pagebreak

\section*{Abkürzungsverzeichnis}
\begin{acronym}
\acro{PMS}{Performance Management System}
\end{acronym}
\pagebreak

% Generate a table of contents?
$GENERATE_TOC$

% BEGIN BODY
\pagenumbering{$PAGE_NUMBERING$}
\setcounter{page}{1}
\setcounter{table}{0}
$body$

\pagenumbering{Roman}
\setcounter{page}{6}
\pagebreak
\listoftables
\pagebreak
\listoffigures
\end{document} 
 

% Include document settings
%\input{settings/document}

% Load language specific Strings
%%!TEX root = ../main.tex

\newcommand{\degreePhrase}{for the}
\newcommand{\departmentPhrase}{from the Course of Studies}
\newcommand{\locationUniversityPhrase}{at the Cooperative State University Baden-Württemberg}
\newcommand{\documentAuthorPhrase}{by}
\newcommand{\documentPeriodPhrase}{Time of Project}
\newcommand{\matriculationNumberPhrase}{Student ID}
\newcommand{\coursePhrase}{Course}
\newcommand{\companyPhrase}{Company}

\newcommand{\tutorPhrase}{%
	\IfStrEqCase{\documentType}{%
		    {T2\_1000}{Reviewer}%
      	{T2\_2000}{Reviewer}%
        {T2\_3000}{Reviewer}%
      	{T2\_3100}{Reviewer}%
      	{T2\_3300}{First Reviewer}%
	}
}

\newcommand{\evaluatorPhrase}{Second Reviewer}
\newcommand{\restrictionNoticesPhrase}{Confidentiality Statement} %Restriction Notice}
\newcommand{\declarationPhrase}{Author's declaration}
\newcommand{\glossaryPhrase}{Glossary}
\newcommand{\langabkverz}{Acronyms}
\newcommand{\langabstract}{Abstract}
\newcommand{\appendixPhrase}{Appendix}
\newcommand{\listingPhrase}{Listing}
\newcommand{\listListingPhrase}{Listings}
\newcommand{\authorListingPhrase}{Listing}
\newcommand{\acronymsPhrase}{Acronyms}

\newcommand{\documentTypePhrase}{%
  \IfStrEqCase{\documentType}{%
      {T2\_1000}{Project Thesis}%
      {T2\_2000}{Project Thesis}%
      {T2\_3000}{Project Thesis}%
      {T2\_3100}{Semester Paper}%
      {T2\_3300}{Bachelor Thesis}%
  }
}

% Add comment feature
\newcommand{\comment}[1]{\par {\bfseries \color{blue} #1 \par}}


%%%%%%% Package Includes %%%%%%%
\usepackage{chngcntr}
\usepackage{float}
\usepackage{subfig}
\usepackage{tabularx}
%\usepackage[margin=\margin,foot=1cm]{geometry}
%\usepackage[activate]{microtype}                                         
%\usepackage[onehalfspacing]{setspace}
\usepackage{makeidx}
\usepackage[autostyle=true,german=quotes]{csquotes}
\usepackage{longtable}
\usepackage{enumitem}	                                                 
\usepackage{graphicx}
\usepackage{pdfpages}                                                         
\usepackage{xcolor} 	                                                    
\usepackage{float}
\usepackage{array}
\usepackage{calc}		                         
\usepackage[right]{eurosym}
\usepackage{wrapfig}
\usepackage{pgffor}                                                             
\usepackage[perpage, hang, multiple, stable]{footmisc}  
\usepackage[printonlyused]{acronym}                                 
\usepackage{listings}
\usepackage[obeyFinal,backgroundcolor=yellow,linecolor=black]{todonotes}
\usepackage{rotating}
\usepackage{lscape}
\usepackage{amsmath}
\usepackage{amssymb}
\usepackage{eurosym}
\usepackage{tocbibind} %makes list of figures, list of tables etc. show up in toc
%\usepackage{\documentFont}
%\usepackage[
%	pdftitle={\documentTitle},
%	pdfauthor={\documentAuthor},
%	pdfsubject={\documentType},
%	pdfcreator={pdflatex, LaTeX with KOMA-Script},
%	pdfpagemode=UseOutlines,       % Show Contents while opening
	%pdfdisplaydoctitle=true, 		% Show document title instead of file name
	%pdflang={\documentLanguage}, % Document language
%]{hyperref}
\usepackage{bookmark}
\usepackage[nonumberlist,toc,xindy]{glossaries}
% add first occurence in footnote
\defglsdisplayfirst[main]{#1#4\protect\footnote{#2}}

% Commands for check and x marks
\usepackage{amssymb}% http://ctan.org/pkg/amssymb
\usepackage{pifont}% http://ctan.org/pkg/pifont
\newcommand{\cmark}{\ding{51}}%
\newcommand{\xmark}{\ding{55}}%
    
% Generate glossary
\makeglossaries{}

% Load colors
%\defineColors{}

% Set Titel, Autor and Date
%\title{\documentTitle}
%\author{\documentAuthor}
%\date{\datum}


% PDF link settings
\hypersetup{%
	colorlinks=true, 		
	linkcolor=LinkColor, 	
	citecolor=LinkColor,
	filecolor=LinkColor,
	menucolor=LinkColor,
	urlcolor=LinkColor,
	linktocpage=true, 
	bookmarksnumbered=true 
}

% Captions fontsize
\addtokomafont{caption}{\small}

% Hurenkinder und Schusterjungen verhindern
% http://projekte.dante.de/DanteFAQ/Silbentrennung
\clubpenalty = 10000 % schließt Schusterjungen aus (Seitenumbruch nach der ersten Zeile eines neuen Absatzes)
\widowpenalty = 10000 % schließt Hurenkinder aus (die letzte Zeile eines Absatzes steht auf einer neuen Seite)
\displaywidowpenalty=10000

% Graphicspath
\graphicspath{{images/}}

% frequently used programing languages
\lstloadlanguages{PHP,Python,Java,C,C++,bash}

%\listingsettings{}
% Rename Listings
\renewcommand\lstlistingname{\listingPhrase}
\renewcommand\lstlistlistingname{\listListingPhrase}
\def\lstlistingautorefname{\authorListingPhrase}

% Spaces in tables
%\setlength{\tabcolsep}{\tableColumnMargin}
%\renewcommand{\arraystretch}{\tableRowMargin}

% Manage ToDo's
%\usepackage{todonotes}
%\usepackage{ulem}
\iflang{de}{\usepackage[english, ngerman]{babel}}
\iflang{en}{\usepackage[ngerman, english]{babel}}
%%%%HEADER%%%%

% The following packages are needed for the combined use of Pandoc and XeTeX engine
\usepackage{fontspec,xltxtra,xunicode}

% These directives are necessary for Pandoc to correctly set the language of the
% document. If the user uses XeLaTeX (the default for Zettlr), it will use
% polyglossia, else it will use babel. We've simply copied that from the default
% template -- thanks to the Pandoc community at this point!
$if(lang)$
\ifxetex
  \usepackage{polyglossia}
  \setmainlanguage[$polyglossia-lang.options$]{$polyglossia-lang.name$}
$for(polyglossia-otherlangs)$
  \setotherlanguage[$polyglossia-otherlangs.options$]{$polyglossia-otherlangs.name$}
$endfor$
\else
  \usepackage[shorthands=off,$for(babel-otherlangs)$$babel-otherlangs$,$endfor$main=$babel-lang$]{babel}
\fi
$endif$

% Captions for tables and images
\usepackage{caption}
\captionsetup{format=plain, font=small, labelfont=bf}

% This package is needed for better rendering of images (including captions)
\usepackage{graphicx}

% These packages are needed by pandoc for certain special glyphs, such as the
% \square symbol for checkboxes.
\usepackage{amssymb,amsmath}

% Calc is needed to compute the image width (to prevent up-scaling of small images)
\usepackage{calc}

% ulem is needed for underlining and strikethrough text commands.
% the option "normalem" is needed to preserve _italics_. Without the option,
% _italics_ will become underlines, which is not desirable.
\usepackage[normalem]{ulem}

% Image scaling code from pandoc default.latex template
\makeatletter
\def\maxwidth{\ifdim\Gin@nat@width>\linewidth\linewidth\else\Gin@nat@width\fi}
\def\maxheight{\ifdim\Gin@nat@height>\textheight\textheight\else\Gin@nat@height\fi}
\makeatother
% Scale images if necessary, so that they will not overflow the page
% margins by default, and it is still possible to overwrite the defaults
% using explicit options in \includegraphics[width, height, ...]{}
\setkeys{Gin}{width=\maxwidth,height=\maxheight,keepaspectratio}

% Now we also have to make sure that the figures Pandoc inserts are positioned
% approximately where they are at. Therefore we have to overwrite the
% \begin{figure} command a little bit by using the float package:
\usepackage{float}
\floatplacement{figure}{H} % ensure default position as H (exactly where the figure is in the text)

% Needed for pandoc's table generation
\usepackage{longtable, booktabs, array}

% Optional packages that Zettlr makes use of
% fontspec is used to set the different fonts to the user's whishes
\setmainfont{$MAIN_FONT$}
\setsansfont{$SANS_FONT$}

% With geometry we can let the user decide on the margin
\usepackage{geometry}
% Additional options for geometry can be found here: http://texdoc.net/texmf-dist/doc/latex/geometry/geometry.pdf
\geometry{verbose,$PAPER_TYPE$,tmargin=$TOP_MARGIN$,bmargin=$BOTTOM_MARGIN$,lmargin=$LEFT_MARGIN$,rmargin=$RIGHT_MARGIN$}

% PAGE NUMBERING OPTIONS: https://de.sharelatex.com/learn/Page_numbering
\pagenumbering{$PAGE_NUMBERING$}

% Line spacing
\usepackage{setspace}

% Can be a multiple of 1 (i.e. 1.5 for 150 percent)
\setstretch{$LINE_SPACING$}

% PDF metadata generated by Zettlr or given by the user, also some other URL
% options (breaklinks breaks URLs so they aren't ugly)
%\usepackage[hyphens]{url} % This package breaks links even better
\usepackage{hyperref}
\hypersetup{
unicode=true,
breaklinks=true,
colorlinks=true,
linkcolor=[rgb]{0.00,0.00,0.00},
urlcolor=[rgb]{0.00,0.00,0.00},
pdftitle={$PDF_TITLE$},
pdfsubject={$PDF_SUBJECT$},
pdfauthor={$PDF_AUTHOR$},
pdfkeywords={$PDF_KEYWORDS$},
pdfproducer={Zettlr with Pandoc and XeLaTeX},
pdfborder={0 0 0}
}

% Setup pandoc code highlighting
% pandoc uses an if statement, but until we have code
% to enable or disable highlighting we'll ignore that.

$highlighting-macros$


% Set the date (either \today or a date from the frontmatter)
\date{$PDF_DATE$}

% The following commands must be provided for Pandoc to work correctly.
\providecommand{\tightlist}{%
  \setlength{\itemsep}{0pt}\setlength{\parskip}{0pt}}

% The following commands are used by Pandoc if a CSL bibliography is present
% to create proper indentation, especially for styles with hanging indentation.
$if(csl-refs)$
\newlength{\cslhangindent}
\setlength{\cslhangindent}{1.5em}
\newlength{\csllabelwidth}
\setlength{\csllabelwidth}{3em}
\newenvironment{CSLReferences}[2] % #1 hanging-ident, #2 entry spacing
 {% don't indent paragraphs
  \setlength{\parindent}{0pt}
  % turn on hanging indent if param 1 is 1
  \ifodd #1 \everypar{\setlength{\hangindent}{\cslhangindent}}\ignorespaces\fi
  % set entry spacing
  \ifnum #2 > 0
  \setlength{\parskip}{#2\baselineskip}
  \fi
 }%
 {}
\usepackage{calc}
\newcommand{\CSLBlock}[1]{#1\hfill\break}
\newcommand{\CSLLeftMargin}[1]{\parbox[t]{\csllabelwidth}{#1}}
\newcommand{\CSLRightInline}[1]{\parbox[t]{\linewidth - \csllabelwidth}{#1}\break}
\newcommand{\CSLIndent}[1]{\hspace{\cslhangindent}#1}
$endif$

\title{$PDF_TITLE$}
\author{$PDF_AUTHOR$}

%\input{/home/damion/Desktop/Studium/Studienarbeit/Export/settings/document.tex}
%\input{/home/damion/Desktop/Studium/Studienarbeit/Export/ads/header.tex}

\begin{document}


% Title page?
%$TITLEPAGE$
\pagenumbering{Roman}
%!TEX root = ../main.tex

\begin{titlepage}
	\begin{longtable}{p{8cm} p{8cm}}
		\raggedright {\raisebox{\ht\strutbox-\totalheight}{\includegraphics[height=2.5cm]{images/cover/logo-dhbw.png}}} &
		\ifDocType{T2\_3100}{%
			% no company logo for semester paper
		}{%
			\raggedleft {\raisebox{\ht\strutbox-\totalheight}{\includegraphics[height=2.5cm]{images/cover/logo-company.png}}}
		}%
	\end{longtable}
	\enlargethispage{20mm}
	\begin{center}
	\doublespacing{
		\vspace*{12mm}	{\LARGE\textbf \documentTitle }}\\
		\vspace*{12mm}	{\large\textbf {\documentTypePhrase}}\\
		
		% degre only for bachelor thesis
		\ifDocType{T2\_3300}{
			\vspace*{12mm}	\degreePhrase\\
			\vspace*{3mm}		{\textbf \degree}\\
		}

		\vspace*{12mm}	\departmentPhrase{} \department\\
	\vspace*{0mm}		\locationUniversityPhrase{} \locationUniversity\\
		\vspace*{12mm}	\documentAuthorPhrase\\
		\vspace*{3mm}		{\large\textbf \documentAuthor}\\
		\vspace*{12mm}	\releaseDate\\
	\end{center}
	\vfill
	\begin{spacing}{1.2}
	\begin{tabbing}
		mmmmmmmmmmmmmmmmmmmmmmmmmm             \= \kill
		\textbf{\documentPeriodPhrase}       \>  \documentPeriod\\
		\textbf{\matriculationNumberPhrase, \coursePhrase}  \>  \matriculationNumber, \course\\

		\ifDocType{T2\_3100}{%
			% no company for the semester paper
		}{%
			\textbf{\companyPhrase}                  \>  \companyName, \companyLocation\\
		}%
		\textbf{\tutorPhrase}               \>  \tutor\\

		% evaluator only for bachelor thesis
		\ifDocType{T2\_3300}{%
			\textbf{\evaluatorPhrase}              \>  \evaluator
		}{}
	\end{tabbing}
	\end{spacing}
\end{titlepage}

%\thispagestyle{empty}

\section*{Erklärung}

\vspace*{2em}

gemäß Ziffer 1.1.13 der Anlage 1 zu §§ 3, 4 und 5  der Studien- und Prüfungsordnung für die Bachelorstudiengänge im Studienbereich Technik der Dualen Hochschule Baden-Württemberg vom 29.09.2017.
  
\ifMultipleAuthors{Wir versichern}{Ich versichere} hiermit, dass \ifMultipleAuthors{wir unsere}{ich meine} 
{\documentTypePhrase} mit dem Thema: 
\vspace{3em}
\begin{addmargin}[25pt]{0pt}
\Large {\itshape \documentTitle }
\end{addmargin}
\vspace{3em}
selbstständig verfasst und  keine anderen als die angegebenen Quellen und Hilfsmittel benutzt \ifMultipleAuthors{haben}{habe}. 
\ifMultipleAuthors{Wir versichern}{Ich versichere} zudem, dass die eingereichte elektronische Fassung mit der gedruckten Fassung 
übereinstimmt. 

\vspace{3em}

\releaseLocation, \releaseDate
\vspace{4em}

\rule{6cm}{0.4pt}

\documentAuthor
\pagebreak

%!TEX root = ../main.tex

%\thispagestyle{empty}

\section*{Sperrvermerk}

\vspace*{2em}

Der Inhalt dieser Arbeit darf weder als Ganzes noch in Auszügen Personen außerhalb des Prüfungsprozesses und des Evaluationsverfahrens zugänglich gemacht werden, sofern keine anderslautende Genehmigung vom Dualen Partner vorliegt.

\vspace{3em}

\releaseLocation, \releaseDate
\vspace{4em}

\rule{6cm}{0.4pt}

\documentAuthor
\pagebreak

\section*{Abkürzungsverzeichnis}
\begin{acronym}
\acro{PMS}{Performance Management System}
\end{acronym}
\pagebreak

% Generate a table of contents?
$GENERATE_TOC$

% BEGIN BODY
\pagenumbering{$PAGE_NUMBERING$}
\setcounter{page}{1}
\setcounter{table}{0}
$body$

\pagenumbering{Roman}
\setcounter{page}{6}
\pagebreak
\listoftables
\pagebreak
\listoffigures
\end{document} 
 

% Include document settings
%\input{settings/document}

% Load language specific Strings
%%!TEX root = ../main.tex

\newcommand{\degreePhrase}{for the}
\newcommand{\departmentPhrase}{from the Course of Studies}
\newcommand{\locationUniversityPhrase}{at the Cooperative State University Baden-Württemberg}
\newcommand{\documentAuthorPhrase}{by}
\newcommand{\documentPeriodPhrase}{Time of Project}
\newcommand{\matriculationNumberPhrase}{Student ID}
\newcommand{\coursePhrase}{Course}
\newcommand{\companyPhrase}{Company}

\newcommand{\tutorPhrase}{%
	\IfStrEqCase{\documentType}{%
		    {T2\_1000}{Reviewer}%
      	{T2\_2000}{Reviewer}%
        {T2\_3000}{Reviewer}%
      	{T2\_3100}{Reviewer}%
      	{T2\_3300}{First Reviewer}%
	}
}

\newcommand{\evaluatorPhrase}{Second Reviewer}
\newcommand{\restrictionNoticesPhrase}{Confidentiality Statement} %Restriction Notice}
\newcommand{\declarationPhrase}{Author's declaration}
\newcommand{\glossaryPhrase}{Glossary}
\newcommand{\langabkverz}{Acronyms}
\newcommand{\langabstract}{Abstract}
\newcommand{\appendixPhrase}{Appendix}
\newcommand{\listingPhrase}{Listing}
\newcommand{\listListingPhrase}{Listings}
\newcommand{\authorListingPhrase}{Listing}
\newcommand{\acronymsPhrase}{Acronyms}

\newcommand{\documentTypePhrase}{%
  \IfStrEqCase{\documentType}{%
      {T2\_1000}{Project Thesis}%
      {T2\_2000}{Project Thesis}%
      {T2\_3000}{Project Thesis}%
      {T2\_3100}{Semester Paper}%
      {T2\_3300}{Bachelor Thesis}%
  }
}

% Add comment feature
\newcommand{\comment}[1]{\par {\bfseries \color{blue} #1 \par}}


%%%%%%% Package Includes %%%%%%%
\usepackage{chngcntr}
\usepackage{float}
\usepackage{subfig}
\usepackage{tabularx}
%\usepackage[margin=\margin,foot=1cm]{geometry}
%\usepackage[activate]{microtype}                                         
%\usepackage[onehalfspacing]{setspace}
\usepackage{makeidx}
\usepackage[autostyle=true,german=quotes]{csquotes}
\usepackage{longtable}
\usepackage{enumitem}	                                                 
\usepackage{graphicx}
\usepackage{pdfpages}                                                         
\usepackage{xcolor} 	                                                    
\usepackage{float}
\usepackage{array}
\usepackage{calc}		                         
\usepackage[right]{eurosym}
\usepackage{wrapfig}
\usepackage{pgffor}                                                             
\usepackage[perpage, hang, multiple, stable]{footmisc}  
\usepackage[printonlyused]{acronym}                                 
\usepackage{listings}
\usepackage[obeyFinal,backgroundcolor=yellow,linecolor=black]{todonotes}
\usepackage{rotating}
\usepackage{lscape}
\usepackage{amsmath}
\usepackage{amssymb}
\usepackage{eurosym}
\usepackage{tocbibind} %makes list of figures, list of tables etc. show up in toc
%\usepackage{\documentFont}
%\usepackage[
%	pdftitle={\documentTitle},
%	pdfauthor={\documentAuthor},
%	pdfsubject={\documentType},
%	pdfcreator={pdflatex, LaTeX with KOMA-Script},
%	pdfpagemode=UseOutlines,       % Show Contents while opening
	%pdfdisplaydoctitle=true, 		% Show document title instead of file name
	%pdflang={\documentLanguage}, % Document language
%]{hyperref}
\usepackage{bookmark}
\usepackage[nonumberlist,toc,xindy]{glossaries}
% add first occurence in footnote
\defglsdisplayfirst[main]{#1#4\protect\footnote{#2}}

% Commands for check and x marks
\usepackage{amssymb}% http://ctan.org/pkg/amssymb
\usepackage{pifont}% http://ctan.org/pkg/pifont
\newcommand{\cmark}{\ding{51}}%
\newcommand{\xmark}{\ding{55}}%
    
% Generate glossary
\makeglossaries{}

% Load colors
%\defineColors{}

% Set Titel, Autor and Date
%\title{\documentTitle}
%\author{\documentAuthor}
%\date{\datum}


% PDF link settings
\hypersetup{%
	colorlinks=true, 		
	linkcolor=LinkColor, 	
	citecolor=LinkColor,
	filecolor=LinkColor,
	menucolor=LinkColor,
	urlcolor=LinkColor,
	linktocpage=true, 
	bookmarksnumbered=true 
}

% Captions fontsize
\addtokomafont{caption}{\small}

% Hurenkinder und Schusterjungen verhindern
% http://projekte.dante.de/DanteFAQ/Silbentrennung
\clubpenalty = 10000 % schließt Schusterjungen aus (Seitenumbruch nach der ersten Zeile eines neuen Absatzes)
\widowpenalty = 10000 % schließt Hurenkinder aus (die letzte Zeile eines Absatzes steht auf einer neuen Seite)
\displaywidowpenalty=10000

% Graphicspath
\graphicspath{{images/}}

% frequently used programing languages
\lstloadlanguages{PHP,Python,Java,C,C++,bash}

%\listingsettings{}
% Rename Listings
\renewcommand\lstlistingname{\listingPhrase}
\renewcommand\lstlistlistingname{\listListingPhrase}
\def\lstlistingautorefname{\authorListingPhrase}

% Spaces in tables
%\setlength{\tabcolsep}{\tableColumnMargin}
%\renewcommand{\arraystretch}{\tableRowMargin}

% Manage ToDo's
%\usepackage{todonotes}
%\usepackage{ulem}
\iflang{de}{\usepackage[english, ngerman]{babel}}
\iflang{en}{\usepackage[ngerman, english]{babel}}
%%%%HEADER%%%%

% The following packages are needed for the combined use of Pandoc and XeTeX engine
\usepackage{fontspec,xltxtra,xunicode}

% These directives are necessary for Pandoc to correctly set the language of the
% document. If the user uses XeLaTeX (the default for Zettlr), it will use
% polyglossia, else it will use babel. We've simply copied that from the default
% template -- thanks to the Pandoc community at this point!
$if(lang)$
\ifxetex
  \usepackage{polyglossia}
  \setmainlanguage[$polyglossia-lang.options$]{$polyglossia-lang.name$}
$for(polyglossia-otherlangs)$
  \setotherlanguage[$polyglossia-otherlangs.options$]{$polyglossia-otherlangs.name$}
$endfor$
\else
  \usepackage[shorthands=off,$for(babel-otherlangs)$$babel-otherlangs$,$endfor$main=$babel-lang$]{babel}
\fi
$endif$

% Captions for tables and images
\usepackage{caption}
\captionsetup{format=plain, font=small, labelfont=bf}

% This package is needed for better rendering of images (including captions)
\usepackage{graphicx}

% These packages are needed by pandoc for certain special glyphs, such as the
% \square symbol for checkboxes.
\usepackage{amssymb,amsmath}

% Calc is needed to compute the image width (to prevent up-scaling of small images)
\usepackage{calc}

% ulem is needed for underlining and strikethrough text commands.
% the option "normalem" is needed to preserve _italics_. Without the option,
% _italics_ will become underlines, which is not desirable.
\usepackage[normalem]{ulem}

% Image scaling code from pandoc default.latex template
\makeatletter
\def\maxwidth{\ifdim\Gin@nat@width>\linewidth\linewidth\else\Gin@nat@width\fi}
\def\maxheight{\ifdim\Gin@nat@height>\textheight\textheight\else\Gin@nat@height\fi}
\makeatother
% Scale images if necessary, so that they will not overflow the page
% margins by default, and it is still possible to overwrite the defaults
% using explicit options in \includegraphics[width, height, ...]{}
\setkeys{Gin}{width=\maxwidth,height=\maxheight,keepaspectratio}

% Now we also have to make sure that the figures Pandoc inserts are positioned
% approximately where they are at. Therefore we have to overwrite the
% \begin{figure} command a little bit by using the float package:
\usepackage{float}
\floatplacement{figure}{H} % ensure default position as H (exactly where the figure is in the text)

% Needed for pandoc's table generation
\usepackage{longtable, booktabs, array}

% Optional packages that Zettlr makes use of
% fontspec is used to set the different fonts to the user's whishes
\setmainfont{$MAIN_FONT$}
\setsansfont{$SANS_FONT$}

% With geometry we can let the user decide on the margin
\usepackage{geometry}
% Additional options for geometry can be found here: http://texdoc.net/texmf-dist/doc/latex/geometry/geometry.pdf
\geometry{verbose,$PAPER_TYPE$,tmargin=$TOP_MARGIN$,bmargin=$BOTTOM_MARGIN$,lmargin=$LEFT_MARGIN$,rmargin=$RIGHT_MARGIN$}

% PAGE NUMBERING OPTIONS: https://de.sharelatex.com/learn/Page_numbering
\pagenumbering{$PAGE_NUMBERING$}

% Line spacing
\usepackage{setspace}

% Can be a multiple of 1 (i.e. 1.5 for 150 percent)
\setstretch{$LINE_SPACING$}

% PDF metadata generated by Zettlr or given by the user, also some other URL
% options (breaklinks breaks URLs so they aren't ugly)
%\usepackage[hyphens]{url} % This package breaks links even better
\usepackage{hyperref}
\hypersetup{
unicode=true,
breaklinks=true,
colorlinks=true,
linkcolor=[rgb]{0.00,0.00,0.00},
urlcolor=[rgb]{0.00,0.00,0.00},
pdftitle={$PDF_TITLE$},
pdfsubject={$PDF_SUBJECT$},
pdfauthor={$PDF_AUTHOR$},
pdfkeywords={$PDF_KEYWORDS$},
pdfproducer={Zettlr with Pandoc and XeLaTeX},
pdfborder={0 0 0}
}

% Setup pandoc code highlighting
% pandoc uses an if statement, but until we have code
% to enable or disable highlighting we'll ignore that.

$highlighting-macros$


% Set the date (either \today or a date from the frontmatter)
\date{$PDF_DATE$}

% The following commands must be provided for Pandoc to work correctly.
\providecommand{\tightlist}{%
  \setlength{\itemsep}{0pt}\setlength{\parskip}{0pt}}

% The following commands are used by Pandoc if a CSL bibliography is present
% to create proper indentation, especially for styles with hanging indentation.
$if(csl-refs)$
\newlength{\cslhangindent}
\setlength{\cslhangindent}{1.5em}
\newlength{\csllabelwidth}
\setlength{\csllabelwidth}{3em}
\newenvironment{CSLReferences}[2] % #1 hanging-ident, #2 entry spacing
 {% don't indent paragraphs
  \setlength{\parindent}{0pt}
  % turn on hanging indent if param 1 is 1
  \ifodd #1 \everypar{\setlength{\hangindent}{\cslhangindent}}\ignorespaces\fi
  % set entry spacing
  \ifnum #2 > 0
  \setlength{\parskip}{#2\baselineskip}
  \fi
 }%
 {}
\usepackage{calc}
\newcommand{\CSLBlock}[1]{#1\hfill\break}
\newcommand{\CSLLeftMargin}[1]{\parbox[t]{\csllabelwidth}{#1}}
\newcommand{\CSLRightInline}[1]{\parbox[t]{\linewidth - \csllabelwidth}{#1}\break}
\newcommand{\CSLIndent}[1]{\hspace{\cslhangindent}#1}
$endif$

\title{$PDF_TITLE$}
\author{$PDF_AUTHOR$}

%\input{/home/damion/Desktop/Studium/Studienarbeit/Export/settings/document.tex}
%\input{/home/damion/Desktop/Studium/Studienarbeit/Export/ads/header.tex}

\begin{document}


% Title page?
%$TITLEPAGE$
\pagenumbering{Roman}
%!TEX root = ../main.tex

\begin{titlepage}
	\begin{longtable}{p{8cm} p{8cm}}
		\raggedright {\raisebox{\ht\strutbox-\totalheight}{\includegraphics[height=2.5cm]{images/cover/logo-dhbw.png}}} &
		\ifDocType{T2\_3100}{%
			% no company logo for semester paper
		}{%
			\raggedleft {\raisebox{\ht\strutbox-\totalheight}{\includegraphics[height=2.5cm]{images/cover/logo-company.png}}}
		}%
	\end{longtable}
	\enlargethispage{20mm}
	\begin{center}
	\doublespacing{
		\vspace*{12mm}	{\LARGE\textbf \documentTitle }}\\
		\vspace*{12mm}	{\large\textbf {\documentTypePhrase}}\\
		
		% degre only for bachelor thesis
		\ifDocType{T2\_3300}{
			\vspace*{12mm}	\degreePhrase\\
			\vspace*{3mm}		{\textbf \degree}\\
		}

		\vspace*{12mm}	\departmentPhrase{} \department\\
	\vspace*{0mm}		\locationUniversityPhrase{} \locationUniversity\\
		\vspace*{12mm}	\documentAuthorPhrase\\
		\vspace*{3mm}		{\large\textbf \documentAuthor}\\
		\vspace*{12mm}	\releaseDate\\
	\end{center}
	\vfill
	\begin{spacing}{1.2}
	\begin{tabbing}
		mmmmmmmmmmmmmmmmmmmmmmmmmm             \= \kill
		\textbf{\documentPeriodPhrase}       \>  \documentPeriod\\
		\textbf{\matriculationNumberPhrase, \coursePhrase}  \>  \matriculationNumber, \course\\

		\ifDocType{T2\_3100}{%
			% no company for the semester paper
		}{%
			\textbf{\companyPhrase}                  \>  \companyName, \companyLocation\\
		}%
		\textbf{\tutorPhrase}               \>  \tutor\\

		% evaluator only for bachelor thesis
		\ifDocType{T2\_3300}{%
			\textbf{\evaluatorPhrase}              \>  \evaluator
		}{}
	\end{tabbing}
	\end{spacing}
\end{titlepage}

%\thispagestyle{empty}

\section*{Erklärung}

\vspace*{2em}

gemäß Ziffer 1.1.13 der Anlage 1 zu §§ 3, 4 und 5  der Studien- und Prüfungsordnung für die Bachelorstudiengänge im Studienbereich Technik der Dualen Hochschule Baden-Württemberg vom 29.09.2017.
  
\ifMultipleAuthors{Wir versichern}{Ich versichere} hiermit, dass \ifMultipleAuthors{wir unsere}{ich meine} 
{\documentTypePhrase} mit dem Thema: 
\vspace{3em}
\begin{addmargin}[25pt]{0pt}
\Large {\itshape \documentTitle }
\end{addmargin}
\vspace{3em}
selbstständig verfasst und  keine anderen als die angegebenen Quellen und Hilfsmittel benutzt \ifMultipleAuthors{haben}{habe}. 
\ifMultipleAuthors{Wir versichern}{Ich versichere} zudem, dass die eingereichte elektronische Fassung mit der gedruckten Fassung 
übereinstimmt. 

\vspace{3em}

\releaseLocation, \releaseDate
\vspace{4em}

\rule{6cm}{0.4pt}

\documentAuthor
\pagebreak

%!TEX root = ../main.tex

%\thispagestyle{empty}

\section*{Sperrvermerk}

\vspace*{2em}

Der Inhalt dieser Arbeit darf weder als Ganzes noch in Auszügen Personen außerhalb des Prüfungsprozesses und des Evaluationsverfahrens zugänglich gemacht werden, sofern keine anderslautende Genehmigung vom Dualen Partner vorliegt.

\vspace{3em}

\releaseLocation, \releaseDate
\vspace{4em}

\rule{6cm}{0.4pt}

\documentAuthor
\pagebreak

\section*{Abkürzungsverzeichnis}
\begin{acronym}
\acro{PMS}{Performance Management System}
\end{acronym}
\pagebreak

% Generate a table of contents?
$GENERATE_TOC$

% BEGIN BODY
\pagenumbering{$PAGE_NUMBERING$}
\setcounter{page}{1}
\setcounter{table}{0}
$body$

\pagenumbering{Roman}
\setcounter{page}{6}
\pagebreak
\listoftables
\pagebreak
\listoffigures
\end{document} 
